\documentclass[11pt]{article}

\usepackage[letterpaper,margin=0.9in]{geometry}
\usepackage[T1]{fontenc}
\usepackage[charter]{mathdesign}
\usepackage[scaled=0.92]{helvet}
\usepackage{microtype}
\usepackage{amsmath,amsthm}
\usepackage{bm}
\usepackage{enumitem}
\usepackage{xcolor}
\usepackage{titlesec}
\usepackage[hidelinks]{hyperref}

\definecolor{headblue}{RGB}{24,54,96}
\setlength{\parindent}{0pt}
\setlength{\parskip}{6pt}
\setlist{nosep,leftmargin=1.4em,topsep=3pt,itemsep=2pt}
\titleformat{\section}{\sffamily\large\bfseries\color{headblue}}{\thesection}{0.6em}{}
\titlespacing*{\section}{0pt}{14pt}{5pt}
\titleformat{\subsection}{\sffamily\normalsize\bfseries}{\thesubsection}{0.5em}{}
\titlespacing*{\subsection}{0pt}{10pt}{3pt}

\theoremstyle{plain}
\newtheorem{prop}{Proposition}
\newtheorem{lem}[prop]{Lemma}
\newtheorem{thm}[prop]{Theorem}
\newtheorem{cor}[prop]{Corollary}
\theoremstyle{remark}
\newtheorem*{rem}{Remark}
\newtheorem*{ex}{Example}

\newcommand{\EIG}{\mathrm{EIG}}
\newcommand{\KL}{\mathrm{KL}}
\newcommand{\Var}{\mathrm{Var}}
\newcommand{\Cov}{\mathrm{Cov}}
\newcommand{\one}{\mathbf{1}}
\newcommand{\lse}{\operatorname{lse}}
\newcommand{\tr}{\operatorname{tr}}

\begin{document}

\begin{center}
{\sffamily\bfseries\LARGE Intelligent Eyes: Derivations}\\[4pt]
{\sffamily\large Companion to the DA411 proposal}\\[8pt]
Dawit Chun
\end{center}
\noindent\rule{\textwidth}{0.4pt}

This note works out every equation used in the proposal, step by step, with its assumptions. Equation numbers (1)--(3) in the proposal correspond to Propositions~\ref{prop:weights}, \ref{prop:pred} and Theorem~\ref{thm:eig} here.

\section{Setup and notation}

\begin{itemize}
  \item Species $y\in\{1,\dots,K\}$; training class frequencies $\pi(y)$.
  \item A clip has $T$ frames. Observation $j=(f,k)$ is frame $f$ seen at position $k$, with $k=0$ the glance (whole frame at 224 px) and $k=1,\dots,9$ the crops of a $3\times3$ grid, each half the frame's side ($2\times$ zoom).
  \item $S$ is the ordered set of observations taken so far (all glances, then the chosen looks). $\one$ is a vector of ones.
  \item A frozen encoder plus a linear classifier gives $q(y\mid f_j)$. The \emph{vote} of observation $j$ is the vector $\ell_j\in\mathbb R^K$ with
  \[ \ell_j(y)=\log q(y\mid f_j)-\log\pi(y). \]
  \item The belief is $p=\operatorname{softmax}(s)$ for a score vector $s\in\mathbb R^K$; $\lse(s)=\log\sum_y e^{s_y}$.
  \item $F=\mathrm{diag}(p)-pp^\top$.
\end{itemize}

\textbf{Assumptions.}
(A1) one species per clip.
(A2) For each species $y$, the votes $(\ell_j(y))_j$ are jointly Gaussian with covariance $\sigma^2R$, where $R$ is a correlation matrix ($R_{jj}=1$); different species components are treated as independent.
(A3) The classifiers are calibrated: $q(y\mid f)=p(y\mid f)$ under the training distribution.
(A4) (Used only for the greedy guarantee.) Observations are conditionally independent given the species.

\section{Votes and Bayes' rule without correlation}

\begin{lem}\label{lem:naive}
Under (A3), if the observations in $S$ are conditionally independent given $y$, then
\[ \log p(y\mid f_S) = \log\pi(y) + \sum_{j\in S}\ell_j(y) + \text{const}. \]
\end{lem}
\begin{proof}
By Bayes' rule and (A3), $q(y\mid f_j)=p(f_j\mid y)\pi(y)/p(f_j)$, so $\ell_j(y)=\log p(f_j\mid y)-\log p(f_j)$. Conditional independence gives $p(f_S\mid y)=\prod_j p(f_j\mid y)$, hence
$\log p(y\mid f_S)=\log\pi(y)+\sum_j\log p(f_j\mid y)-\log p(f_S)=\log\pi(y)+\sum_j\ell_j(y)+\text{const}$,
where the constant $\sum_j\log p(f_j)-\log p(f_S)$ does not depend on $y$.
\end{proof}

So with independent observations the votes simply add. Two overlapping crops, or one spot in two near-identical frames, are not independent, and adding both counts the same evidence twice. The next section fixes this.

\section{\texorpdfstring{Correlated votes: the weights $\bm w=R^{-1}\one$ (proposal eq.~1)}{Correlated votes: the weights (proposal eq. 1)}}

\textbf{Model.} Let $L(y)$ be the evidence that one independent view of the clip would give. Each vote is a noisy copy of it,
\[ \ell_j(y)=L(y)+e_j(y),\qquad \big(e_j(y)\big)_{j\in S}\sim\mathcal N(0,\sigma^2R_S), \]
which is (A2) with a common mean.

\begin{prop}[weights]\label{prop:weights}
Under this model, the best linear unbiased estimate of $L$ from the votes in $S$ is
\[ \hat L=\frac{\one^\top R_S^{-1}\bm\ell_S}{\one^\top R_S^{-1}\one},\qquad \Var(\hat L)=\frac{\sigma^2}{n_{\text{eff}}},\qquad n_{\text{eff}}=\one^\top R_S^{-1}\one . \]
The votes therefore carry the information of $n_{\text{eff}}$ independent views with mean $\hat L$, and Lemma~\ref{lem:naive} applied to those views gives
\[ p(y)\propto\pi(y)\exp\big(n_{\text{eff}}\hat L(y)\big)=\pi(y)\exp\Big(\sum_{j\in S}w_j\ell_j(y)\Big),\qquad \bm w=R_S^{-1}\one . \]
\end{prop}
\begin{proof}
Generalized least squares: minimize $(\bm\ell-L\one)^\top R^{-1}(\bm\ell-L\one)$ over $L$. Setting the derivative $-2\one^\top R^{-1}(\bm\ell-L\one)$ to zero gives $\hat L$. Its variance is $\frac{\one^\top R^{-1}(\sigma^2R)R^{-1}\one}{(\one^\top R^{-1}\one)^2}=\frac{\sigma^2}{\one^\top R^{-1}\one}$. By the Gauss--Markov theorem (applied after whitening by $R^{-1/2}$), no other linear unbiased estimate has smaller variance. Finally, $n_{\text{eff}}\hat L=\one^\top R^{-1}\bm\ell=\sum_j w_j\ell_j$.
\end{proof}

\begin{rem}
With $R=I$ the weights are all 1 and Lemma~\ref{lem:naive} is recovered. Treating correlated votes as $n_{\text{eff}}$ independent views with the same total precision is the modeling step; everything after it is exact under (A2).
\end{rem}

\begin{ex}[two votes]
$R=\begin{pmatrix}1&r\\ r&1\end{pmatrix}$, $R^{-1}=\frac{1}{1-r^2}\begin{pmatrix}1&-r\\-r&1\end{pmatrix}$, so $\bm w=R^{-1}\one=\frac{1-r}{1-r^2}\one=\frac{1}{1+r}\one$ and $n_{\text{eff}}=\frac{2}{1+r}$. Independent votes ($r=0$) count as 2; identical votes ($r\to1$) count as 1.
\end{ex}

\begin{prop}[a repeat cannot change the answer]\label{prop:repeat}
Suppose observation $i\in S$ is fed again as a new observation $j$, so $\ell_j=\ell_i$ and $R$'s new row equals row $i$. Then $R_{S\cup j}$ is singular; with the pseudo-inverse, $\one^\top R_{S\cup j}^{+}\bm\ell_{S\cup j}=\one^\top R_S^{-1}\bm\ell_S$ and $n_{\text{eff}}$ is unchanged.
\end{prop}
\begin{proof}
Let $M=\begin{pmatrix}I\\ e_i^\top\end{pmatrix}$, an $(n{+}1)\times n$ matrix of full column rank. Then $\bm\ell_{S\cup j}=M\bm\ell_S$, $\one_{n+1}=M\one_n$, and $R_{S\cup j}=MR_SM^\top$. For $M$ of full column rank and $R_S$ invertible, $(MR_SM^\top)^+=(M^\top)^+R_S^{-1}M^+$. Hence
$\one^\top M^\top(M^\top)^+R_S^{-1}M^+M\bm\ell_S=\one^\top R_S^{-1}\bm\ell_S$,
because $M^\top(M^\top)^+=I$ ($M^\top$ has full row rank) and $M^+M=I$ ($M$ has full column rank). The same computation with $\bm\ell$ replaced by $\one$ gives the statement for $n_{\text{eff}}$.
\end{proof}

This is the isolation test for the failure mode ``wrong $R$'' in the proposal: feeding a frame twice must leave the prediction unchanged.

\subsection{\texorpdfstring{The structure of $R$}{The structure of R}}
$R_{(f,k),(f',k')}=A_{kk'}\rho^{|f-f'|}$, i.e.\ $R=T_\rho\otimes A$ in frame-major order, where $A$ is the $10\times10$ correlation table of positions (glance and 9 crops) and $(T_\rho)_{ff'}=\rho^{|f-f'|}$ is the correlation matrix of an AR(1) process.
\begin{itemize}
  \item $T_\rho$ is positive definite for $|\rho|<1$, and $A$ is made positive semidefinite by clipping negative eigenvalues and rescaling the diagonal to 1. Eigenvalues of a Kronecker product are products of eigenvalues, so $R$ is positive semidefinite.
  \item \emph{Fitting.} On training clips, take residuals $e_j=\ell_j-\mu_j$ (Section~5). $\hat A_{kk'}$ is the pooled correlation of residuals at positions $k,k'$ in the same frame. $\hat\rho$ minimizes $\sum_{d\ge1}(\hat c_d-\rho^d)^2$, where $\hat c_d$ is the pooled correlation of same-position residuals $d$ frames apart.
\end{itemize}

\begin{prop}[effective number of frames]\label{prop:ar1}
For $T$ glances alone ($R=T_\rho$),
\[ w_1=w_T=\frac{1}{1+\rho},\qquad w_2=\dots=w_{T-1}=\frac{1-\rho}{1+\rho},\qquad n_{\text{eff}}=\frac{T-(T-2)\rho}{1+\rho}. \]
\end{prop}
\begin{proof}
The AR(1) correlation matrix has the tridiagonal inverse $T_\rho^{-1}=\frac{1}{1-\rho^2}\,\mathrm{tridiag}\big(-\rho;\ 1,1+\rho^2,\dots,1+\rho^2,1;\ -\rho\big)$ (check: multiply by $T_\rho$). Row sums give $\bm w$: end rows $\frac{1-\rho}{1-\rho^2}=\frac{1}{1+\rho}$; interior rows $\frac{1+\rho^2-2\rho}{1-\rho^2}=\frac{1-\rho}{1+\rho}$. Summing, $n_{\text{eff}}=\frac{2+(T-2)(1-\rho)}{1+\rho}=\frac{T-(T-2)\rho}{1+\rho}$.
\end{proof}

\begin{ex}
$T=10$, $\rho=0.9$: $n_{\text{eff}}=\frac{10-7.2}{1.9}\approx1.47$. Ten frames of a scene where nothing moves count as about one and a half. With $\rho=0$, $n_{\text{eff}}=T$; as $\rho\to1$, $n_{\text{eff}}\to1$.
\end{ex}

\section{Adding one look: the block inverse}

Let $S$ be the current set and $j$ a candidate. Write
\[ \bm a=R_{SS}^{-1}R_{Sj},\qquad s_j=1-R_{jS}R_{SS}^{-1}R_{Sj}\ \ (\text{the Schur complement}),\qquad \delta_j=\ell_j-R_{jS}R_{SS}^{-1}\bm\ell_S . \]

\begin{prop}[change from one new look]\label{prop:update}
Let $E_S=\one^\top R_{SS}^{-1}\bm\ell_S$ be the total evidence. Adding look $j$ changes it by
\[ \Delta_j = E_{S\cup j}-E_S = c_j\,\delta_j,\qquad c_j=\frac{1-R_{jS}R_{SS}^{-1}\one}{1-R_{jS}R_{SS}^{-1}R_{Sj}}, \]
and the effective count by $n_{\text{eff}}(S\cup j)-n_{\text{eff}}(S)=(1-R_{jS}R_{SS}^{-1}\one)^2/s_j$.
\end{prop}
\begin{proof}
The block inverse is
\[ \begin{pmatrix}R_{SS}&R_{Sj}\\R_{jS}&1\end{pmatrix}^{-1}=\begin{pmatrix}R_{SS}^{-1}+\bm a\bm a^\top/s_j & -\bm a/s_j\\ -\bm a^\top/s_j & 1/s_j\end{pmatrix}. \]
Multiplying by $\one^\top$ on the left and $(\bm\ell_S,\ell_j)$ on the right,
\[ E_{S\cup j}=E_S+\frac{(\one^\top\bm a)(\bm a^\top\bm\ell_S)}{s_j}-\frac{(\one^\top\bm a)\ell_j}{s_j}-\frac{\bm a^\top\bm\ell_S}{s_j}+\frac{\ell_j}{s_j}
=E_S+\frac{(1-\one^\top\bm a)(\ell_j-\bm a^\top\bm\ell_S)}{s_j}. \]
Since $R$ is symmetric, $\one^\top\bm a=R_{jS}R_{SS}^{-1}\one$ and $\bm a^\top\bm\ell_S=R_{jS}R_{SS}^{-1}\bm\ell_S$. Replacing $\bm\ell$ by $\one$ gives the change in $n_{\text{eff}}$.
\end{proof}

\textbf{Meaning.} A new look adds only its \emph{innovation} $\delta_j$: the vote minus what the earlier observations already imply about it. The factor $c_j$ converts it into evidence.

\begin{ex}
One earlier look with correlation $r$: $\bm a=r$, $s_j=1-r^2$, $c_j=\frac{1-r}{1-r^2}=\frac{1}{1+r}$, and $E=\ell_1+\frac{\ell_2-r\ell_1}{1+r}=\frac{\ell_1+\ell_2}{1+r}$, matching the two-vote example. For $r=0.9$: $c_j\approx0.53$.
\end{ex}

The update costs $O(|S|^2)$ with a Cholesky factor of $R_{SS}$, and $|S|\le T+2\le12$.

\section{Predicting a look before taking it (proposal eq.~2)}

\textbf{Ridge regression.} For candidate $j=(f,k)$, let $x_j=[\bar u_{fk};g_f]$ be the mean of the glance's patch tokens inside crop $k$ together with the frame's glance feature. On training looks, with centered $X$ and $Y$ (one row per look, one column per species),
\[ \min_W \|Y-XW\|_F^2+\lambda\|W\|_F^2\ \Longrightarrow\ W=(X^\top X+\lambda I)^{-1}X^\top Y, \]
so $\mu_j=W^\top x_j+b$. The error variance $\sigma^2$ is the mean squared residual on held-out clips.

\begin{prop}[Gaussian conditioning]\label{prop:pred}
Under (A2), with prior predictions $\mu$,
\[ \ell_j\mid\bm\ell_S\ \sim\ \mathcal N\big(\mu_{j\mid S},\ \sigma^2_{j\mid S}\big),\qquad
\mu_{j\mid S}=\mu_j+R_{jS}R_{SS}^{-1}(\bm\ell_S-\bm\mu_S),\qquad
\sigma^2_{j\mid S}=\sigma^2\big(1-R_{jS}R_{SS}^{-1}R_{Sj}\big). \]
\end{prop}
\begin{proof}
The innovation $\delta_j=\ell_j-R_{jS}R_{SS}^{-1}\bm\ell_S$ satisfies $\Cov(\delta_j,\bm\ell_S)=\sigma^2(R_{jS}-R_{jS}R_{SS}^{-1}R_{SS})=0$. Jointly Gaussian and uncorrelated means independent, so $\delta_j$ is independent of $\bm\ell_S$, with mean $\mu_j-R_{jS}R_{SS}^{-1}\bm\mu_S$ and variance $\sigma^2(1-2R_{jS}R_{SS}^{-1}R_{Sj}+R_{jS}R_{SS}^{-1}R_{SS}R_{SS}^{-1}R_{Sj})=\sigma^2 s_j$. Writing $\ell_j=\delta_j+R_{jS}R_{SS}^{-1}\bm\ell_S$ and conditioning on $\bm\ell_S$ gives the result.
\end{proof}

\begin{cor}[distribution of the change]\label{cor:delta}
Combining Propositions~\ref{prop:update} and \ref{prop:pred}, the change in evidence from look $j$ is, for each species component,
\[ \Delta_j\sim\mathcal N(m_j,\ v_j),\qquad m_j=c_j\big(\mu_{j\mid S}-R_{jS}R_{SS}^{-1}\bm\ell_S\big)=c_j\big(\mu_j-R_{jS}R_{SS}^{-1}\bm\mu_S\big),\qquad v_j=c_j^2\,\sigma^2_{j\mid S}. \]
\end{cor}
With no correlation ($R_{jS}=0$), $c_j=1$, $m_j=\mu_j$ and $v_j=\sigma^2$. For $r=0.9$ with one earlier look, $v_j=0.53^2\times0.19\,\sigma^2\approx0.05\,\sigma^2$.

\section{Expected information gain (proposal eq.~3), the result to derive}

\begin{lem}[expected KL is mutual information]\label{lem:mi}
$\displaystyle I(Y;\ell_j)=\mathbb E_{\ell_j}\big[\KL\big(p(\cdot\mid\ell_j)\,\|\,p(\cdot)\big)\big]$.
\end{lem}
\begin{proof}
$I(Y;\ell)=\sum_{\ell}p(\ell)\sum_y p(y\mid\ell)\log\frac{p(y\mid\ell)}{p(y)}$, and the inner sum is the KL divergence (integrals for continuous $\ell$).
\end{proof}

\begin{lem}[the curvature of softmax]\label{lem:hess}
$\nabla\lse(s)=p$ and $\nabla^2\lse(s)=F=\mathrm{diag}(p)-pp^\top$, which is also the Jacobian of $\operatorname{softmax}$.
\end{lem}
\begin{proof}
$\partial_{s_y}\lse=e^{s_y}/\sum_z e^{s_z}=p_y$, and $\partial_{s_z}p_y=p_y(\mathbb 1[y=z]-p_z)$.
\end{proof}

\begin{lem}[second-order KL]\label{lem:taylor}
Let $p=\operatorname{softmax}(s)$ and $p_{\text{new}}=\operatorname{softmax}(s+\Delta)$. Then
\[ \KL(p_{\text{new}}\,\|\,p)=\tfrac12\Delta^\top F\Delta+O(\|\Delta\|^3). \]
\end{lem}
\begin{proof}
$\log p_{\text{new},y}-\log p_y=\Delta_y-[\lse(s+\Delta)-\lse(s)]$, so $\KL=p_{\text{new}}^\top\Delta-[\lse(s+\Delta)-\lse(s)]$. By Lemma~\ref{lem:hess}, $\lse(s+\Delta)-\lse(s)=p^\top\Delta+\frac12\Delta^\top F\Delta+O(\|\Delta\|^3)$ and $p_{\text{new}}=p+F\Delta+O(\|\Delta\|^2)$, so $p_{\text{new}}^\top\Delta=p^\top\Delta+\Delta^\top F\Delta+O(\|\Delta\|^3)$. Subtracting, $\KL=\frac12\Delta^\top F\Delta+O(\|\Delta\|^3)$. (The KL is 0 with zero slope at $\Delta=0$; $F$ is its curvature.)
\end{proof}

\begin{lem}[the variance identity]\label{lem:var}
$\Delta^\top F\Delta=\sum_y p_y\Delta_y^2-\big(\sum_y p_y\Delta_y\big)^2=\Var_{y\sim p}[\Delta(y)]$. Consequently:
\begin{enumerate}[label=(\alph*)]
  \item $F\one=0$: adding the same amount to every species changes nothing, so constants such as $-\log\pi$ shifts and $\lse$ terms drop out;
  \item $F\succeq0$ and $\tr F=1-\sum_y p_y^2$ (the Gini impurity of the belief);
  \item $F\to0$ as $p$ becomes one-hot.
\end{enumerate}
\end{lem}
\begin{proof}
Expand $\Delta^\top(\mathrm{diag}(p)-pp^\top)\Delta$. (a) $F\one=p-p(p^\top\one)=0$. (b) A variance is nonnegative; $\tr F=\sum_yp_y-\sum_yp_y^2$. (c) If $p=e_y$, then $\mathrm{diag}(p)=pp^\top$.
\end{proof}

\begin{thm}[closed-form EIG]\label{thm:eig}
Under (A2) and Corollary~\ref{cor:delta}, to second order in $\Delta$,
\[ \EIG(j)=\mathbb E\big[\KL(p_{\text{new}}\,\|\,p)\big]\approx\tfrac12\big(m_j^\top F\,m_j+v_j\,\tr F\big). \]
\end{thm}
\begin{proof}
By Lemma~\ref{lem:taylor}, $\EIG(j)\approx\frac12\mathbb E[\Delta^\top F\Delta]$. For $\Delta\sim\mathcal N(m,\Sigma)$, $\mathbb E[\Delta^\top F\Delta]=m^\top Fm+\tr(F\Sigma)$ (expand $\Delta=m+\epsilon$; the cross term has mean 0 and $\mathbb E[\epsilon^\top F\epsilon]=\tr(F\,\mathbb E[\epsilon\epsilon^\top])$). With $\Sigma=v_jI$ the result follows.
\end{proof}

\textbf{Reading the two terms.} $m_j^\top Fm_j=\Var_{y\sim p}[m_j(y)]$ is the \emph{expected signal}: large when the look is predicted to vote differently for the species that are still likely. $v_j\tr F$ is the \emph{surprise}: an uncertain prediction may still turn out informative. A repeat has small $c_j$, hence small $m_j$ and $v_j$.

\begin{ex}[two species]
$p=(p,1-p)$ gives $F=p(1-p)\begin{pmatrix}1&-1\\-1&1\end{pmatrix}$, so $m^\top Fm=p(1-p)(m_1-m_2)^2$ and $\tr F=2p(1-p)$. Hence $\EIG\approx\frac12p(1-p)(\Delta\mu^2+2\sigma^2)$ for a fresh look. At $p=0.5$, $\Delta\mu=2$, $\sigma=0.5$: $\frac12\cdot0.25\cdot4.5=0.5625$ nats, below $\ln2\approx0.693$.
\end{ex}

\textbf{Limits of the approximation.} Mutual information is at most $H(p)\le\ln K$, but the quadratic form is unbounded, so strong votes can push it above the ceiling. I therefore (i) report $\min\{\text{closed form},H(p)\}$, (ii) compute the exact EIG by Monte Carlo, drawing $\Delta^{(1)},\dots,\Delta^{(M)}\sim\mathcal N(m_j,v_jI)$ and averaging the exact $\KL(\operatorname{softmax}(s+\Delta^{(i)})\|p)$, with standard error $O(M^{-1/2})$, and (iii) measure how well the two rank looks (proposal metric 5).

\section{The stopping rule}

\begin{cor}\label{cor:stop}
With $p_{\max}=\max_y p_y$, $\EIG(j)\lesssim(1-p_{\max})\big(\|m_j\|^2+v_j\big)$.
\end{cor}
\begin{proof}
$F\succeq0$ gives $m^\top Fm\le\lambda_{\max}(F)\|m\|^2\le\tr F\,\|m\|^2$, and $\tr F=1-\sum_yp_y^2\le1-p_{\max}^2\le2(1-p_{\max})$. Substitute into Theorem~\ref{thm:eig}.
\end{proof}
So once the model is sure, every look is worth almost nothing. Stopping when $\max_j\EIG(j)<\varepsilon$ makes easy clips use fewer passes, and $\varepsilon$ is chosen on the 10\% held-out training clips.

\section{The greedy guarantee}

Let $f(B)=I(Y;\bm\ell_B)$ be the information a set $B$ of looks gives about the species.

\begin{lem}[submodularity]\label{lem:submod}
Under (A4), $f$ is monotone and submodular: for $B\subseteq C$ and $j\notin C$, $0\le f(C\cup j)-f(C)\le f(B\cup j)-f(B)$.
\end{lem}
\begin{proof}
By the chain rule, $f(B\cup j)-f(B)=I(Y;\ell_j\mid\bm\ell_B)=H(\ell_j\mid\bm\ell_B)-H(\ell_j\mid Y,\bm\ell_B)$, which is $\ge0$. Under (A4), $H(\ell_j\mid Y,\bm\ell_B)=H(\ell_j\mid Y)$ does not depend on $B$, so the gain is $H(\ell_j\mid\bm\ell_B)-H(\ell_j\mid Y)$. Conditioning on more cannot increase entropy, so $H(\ell_j\mid\bm\ell_C)\le H(\ell_j\mid\bm\ell_B)$ for $B\subseteq C$.
\end{proof}

\begin{thm}[Nemhauser, Wolsey and Fisher, 1978]\label{thm:greedy}
If $f$ is monotone submodular with $f(\emptyset)=0$, then greedy selection of $k$ elements satisfies
\[ f(G_k)\ \ge\ \big(1-(1-1/k)^k\big)\max_{|B|\le k}f(B)\ \ge\ (1-1/e)\max_{|B|\le k}f(B). \]
For $k=2$ looks the factor is $1-\frac14=\frac34$.
\end{thm}
\begin{proof}
Let $O$ be an optimal set of size $k$ and $G_i$ the greedy set after $i$ steps. By monotonicity and then submodularity,
$f(O)\le f(G_i\cup O)\le f(G_i)+\sum_{o\in O}\big[f(G_i\cup o)-f(G_i)\big]\le f(G_i)+k\big[f(G_{i+1})-f(G_i)\big]$,
because greedy picks the largest single gain. Rearranging, $f(O)-f(G_{i+1})\le(1-\frac1k)\big(f(O)-f(G_i)\big)$. Starting from $f(G_0)=0$ and applying this $k$ times gives $f(O)-f(G_k)\le(1-\frac1k)^kf(O)$. Finally, $(1-\frac1k)^k\le e^{-1}$.
\end{proof}

\textbf{Caveats, tested rather than assumed.} (A4) ignores the correlation that $R$ models, the controller chooses adaptively (look 2 after seeing look 1), and it uses approximate EIG. The bound is therefore a reference point: proposal metric 4 compares the greedy pair with the best of all $\binom{9T}{2}$ pairs.

\section{Numerics and cost}
\begin{itemize}
  \item Heavily overlapping looks make $R_{SS}$ nearly singular. I monitor its condition number and use $R_{SS}+\lambda I$ with small $\lambda$ (Proposition~\ref{prop:repeat} is the limit $\lambda\to0$).
  \item Per decision: $9T\le90$ candidates, each costing $O(|S|^2+K)$ given a Cholesky factor of $R_{SS}$; this is negligible next to one encoder pass.
  \item Encoder passes per clip: $T$ glances plus at most 2 looks, versus $10T$ for zooming everywhere.
\end{itemize}

\section{Statistics used in the evaluation}
\begin{itemize}
  \item \emph{Expected calibration error}, 15 equal-width bins $B_b$ of confidence: $\mathrm{ECE}=\sum_b\frac{|B_b|}{n}\,\big|\mathrm{acc}(B_b)-\mathrm{conf}(B_b)\big|$.
  \item \emph{McNemar test} (paired, same clips): with $b$ clips right only under method 1 and $c$ right only under method 2, the exact two-sided binomial test of $b$ against $\mathrm{Bin}(b+c,\frac12)$.
  \item \emph{Bootstrap intervals}: resample cameras rather than clips, because clips from one camera are correlated, and take the 2.5 and 97.5 percentiles of 10{,}000 resamples.
  \item \emph{Spearman correlation} between predicted and realized gain: $\rho_s=1-\frac{6\sum_id_i^2}{n(n^2-1)}$ with rank differences $d_i$ (average ranks for ties).
\end{itemize}

\section*{References}
{\small
\begin{itemize}[leftmargin=1.2em]
  \item D. V. Lindley, ``On a measure of the information provided by an experiment,'' \emph{Ann. Math. Statist.}, 27(4):986--1005, 1956.
  \item T. M. Cover and J. A. Thomas, \emph{Elements of Information Theory}, 2nd ed., Wiley, 2006 (ch.~2: chain rule, conditioning reduces entropy, mutual information).
  \item G. L. Nemhauser, L. A. Wolsey, and M. L. Fisher, ``An analysis of approximations for maximizing submodular set functions---I,'' \emph{Math. Program.}, 14:265--294, 1978.
  \item A. Krause and C. Guestrin, ``Near-optimal nonmyopic value of information in graphical models,'' in \emph{Proc. UAI}, 2005.
\end{itemize}}

\end{document}
